
\subsection{Published predictors at a glance}
Table~\ref{tab:predictors} collects the methods compared in this paper and
the neighboring literature, grouped by method family. Reported numbers are
the authors' own, on their own datasets; only rows sharing the Veltri 2018
test partition (Table~\ref{tab:sota}) are directly comparable.

\begin{table}[h]\centering\scriptsize
\caption{Published AMP predictors. ``Veltri test'' marks rows evaluated on
the exact Veltri 2018 test partition used here.}
\label{tab:predictors}
\begin{tabular}{p{2.9cm}p{3.1cm}p{4.6cm}p{2.2cm}}
\hline
Method & Family & Signal used & Veltri test \\
\hline
AntiBP2~\cite{antibp2} & SVM, composition & residue and terminal composition & yes (89.37 ACC) \\
iAMP-2L~\cite{iamp2l} & two-level ensemble & pseudo composition, fuzzy KNN & yes (84.90) \\
iAMPpred~\cite{iamppred} & SVM & compositional/physicochemical & yes (88.27) \\
CAMPr3 family~\cite{camp} & ANN/DA/RF/SVM & sequence patterns, signatures & yes (83--88) \\
gkmSVM~\cite{gkm} & string-kernel SVM & gapped $k$-mer spectrum & yes (89.46) \\
AMP Scanner Vr.2~\cite{ampscanner} & CNN & one-hot sequence, DNN & yes (91.29) \\
ACEP~\cite{acep} & CNN + attention & PSSM profiles & yes (93.04) \\
amPEPpy~\cite{ampeppy} & random forest & distribution descriptors & no \\
Macrel~\cite{macrel} & random forest & 22 physicochemical features & no \\
Deep-AmPEP30~\cite{deepampep} & CNN & sequence, 30-aa segments & no \\
PepDesign (this work) & CNN+GNN+RF stack & sequence, chain graph, descriptors & yes (91.22) \\
\hline
\end{tabular}
\end{table}

Three observations position this work. First, the methods sharing the
Veltri partition cluster between 84 and 93\% accuracy, and the top two
(ACEP, AMP Scanner) both use information beyond raw composition:
evolutionary profiles and learned deep features respectively. Our stack
sits between them without profile features. Second, the composition-only
rows (iAMP-2L, iAMPpred, CAMPr3) are 4--9 points behind, consistent with
our own finding that the dipeptide random forest already captures most
compositional signal (Appendix~\ref{app:arch}). Third, the genome-mining
tools (Macrel, amPEPpy) optimize for short-ORF throughput rather than
benchmark accuracy and are therefore not head-to-head competitors; our
screen is closer in spirit to that line, but scores full-length
uncharacterized proteins rather than metagenomic smORFs.
